% =============================================================================
% 第4章：实验设置
% =============================================================================

\section{实验设置}
\label{sec:setup}

本章描述实验的技术配置。数据来源、EDA和预处理流程见第\ref{sec:data}节。

\subsection{损失函数选择：Focal Loss的必要性}

数据集欺诈率仅0.5\%（正负比约1:200），标准的二元交叉熵（BCE）在此极度不平衡条件下存在严重缺陷：正常交易样本占总样本的99.5\%，BCE的梯度将被多数类主导，模型在训练早期就会收敛到"预测所有样本为正常"的退化解。

Focal Loss\cite{lin2017focal}通过两个机制解决这一问题：

\begin{equation}
\text{FL}(p_t) = -\alpha_t (1 - p_t)^{\gamma} \log(p_t)
\label{eq:focal}
\end{equation}

其中$p_t = p$当$y=1$否则$p_t = 1-p$，$\alpha_t = \alpha$当$y=1$否则$\alpha_t = 1-\alpha$。

\begin{itemize}
    \item \textbf{聚焦参数$\gamma=2.0$}：降低易分类样本的损失权重，使训练集中于难分类样本。
    \item \textbf{平衡参数$\alpha=0.25$}：调节正负类权重。本文选择$\alpha=0.25$，因为Sparsemax已通过稀疏特征选择内置了一定的抗不平衡能力。
\end{itemize}


\subsection{模型配置}

本文通过四个控制变量实验构建诊断链，所有变体的架构基础完全一致（表\ref{tab:variant_config}）：

\begin{table}[H]
\centering
\caption{四个实验变体及控制变量。所有变体共享相同的FT-Transformer架构（手写实现，ReGLU，203,841参数），损失函数统一为Focal Loss，LOR-MIM变体额外增加305参数（0.15\%）。}
\label{tab:variant_config}
\begin{tabular}{lccc}
\toprule
\textbf{变体} & \textbf{注意力归一化} & \textbf{数据预处理} \\
\midrule
Softmax基线         & Softmax   & 无 \\
Sparsemax           & Sparsemax & 无 \\
LOR-MIM + Softmax   & Softmax   & LOR-MIM混合 \\
LOR-MIM + Sparsemax & Sparsemax & LOR-MIM混合 \\
\bottomrule
\end{tabular}
\end{table}

共享的架构配置：$d_{model}=64$, $n_{heads}=8$, $n_{layers}=3$, $d_{ffn}=256$，总参数量203,841。

\subsection{LOR-MIM配置}

LOR-MIM的具体原理见第\ref{sec:background}节。其关键超参配置如下：
\begin{itemize}
    \item \textbf{分组数$K=5$}：基于V特征与Amount的互信息经谱聚类将28个V特征分为5组（分组结果：[9, 8, 5, 3, 3]）。
    \item \textbf{门控MLP}：单隐藏层（1$\to$16$\to$5，ReLU激活，Softmax归一化），以标准化的Amount为唯一输入。
    \item \textbf{混合矩阵}：每组一个$d_k \times d_k$矩阵$M_k$，初始化为$I + \mathcal{N}(0, 0.05)$（近单位阵+小高斯噪声）。
    \item \textbf{参数量}：305个额外参数（门控MLP 293 + 混合矩阵12组），占FT-Transformer总参数的0.15\%。
\end{itemize}

\subsection{训练策略}

所有变体使用完全一致的训练配置：AdamW优化器（$lr{=}1{\times}10^{-3}$, weight\_decay$\,{=}\,10^{-5}$, $\beta_1{=}0.9$, $\beta_2{=}0.999$），批量大小512，最大训练轮数100，早停patience=15（基于验证集AUPRC），梯度裁剪max\_norm=1.0，ReduceLROnPlateau学习率调度（patience=8, factor=0.5）。每个模型训练前显式重置random、numpy、torch和cuda的全部随机状态（统一seed=42）以确保可复现性。所有实验在NVIDIA GeForce RTX 4060（8GB）GPU和PyTorch 2.11.0+cu128环境下完成。

\subsection{评估指标}

\begin{itemize}
    \item \textbf{AUPRC}（主要指标）：聚焦正类（欺诈）的排序能力。在不平衡数据上，ROC-AUC的区分度远低于AUPRC\cite{saito2015precision}。
    \item \textbf{AttnStd}（注意力权重标准差）：塌缩程度的直接度量。均匀基线$1/29\approx 0.0345$。
    \item \textbf{ECE}（Expected Calibration Error）：概率校准质量。将预测概率等分为10个分位区间，计算各区间内平均预测概率与平均真实正类率之差的加权平均。ECE越低表示模型的概率估计越可信——这对风控业务至关重要：过度自信的模型会产生误导性高风险预警。
    \item \textbf{F1-Score}、\textbf{Precision}、\textbf{Recall}：在验证集F1最大化的最优阈值处计算。
    \item \textbf{Cost}：业务成本 $= FP \times 1 + FN \times 10$（漏报权重10倍于误报）。
\end{itemize}
